% Generated draft: requires mathematical and source review.
\documentclass[11pt]{article}
\usepackage[margin=25mm]{geometry}
\usepackage{fontspec}
\setmainfont{Latin Modern Roman}
\usepackage{amsmath,amssymb,mathtools,unicode-math}
\setmathfont{Latin Modern Math}
\usepackage{enumitem,needspace,xurl,hyperref}
\setlength{\parindent}{0pt}
\setlength{\parskip}{5pt}
\setlength{\emergencystretch}{3em}
\newcommand{\N}{\mathbb{N}}
\newcommand{\Z}{\mathbb{Z}}
\newcommand{\Q}{\mathbb{Q}}
\newcommand{\R}{\mathbb{R}}
\newcommand{\C}{\mathbb{C}}
\newcommand{\F}{\mathbb{F}}
\newcommand{\A}{\mathbb{A}}
\newcommand{\PP}{\mathbb{P}}
\newcommand{\E}{\mathbb{E}}
\newcommand{\eps}{\varepsilon}
\newcommand{\dd}{\,\mathrm{d}}
\begin{document}
\section*{UnsolvedMath ID: 20000824}
\textbf{Source problem code:} catalogue entry
\subsection{Definitions and mathematical statement}
The source attributes to Joe Harris the question whether a nondegenerate projectively rigid curve in $\PP^n$ can be different from a rational normal curve. Nondegenerate means not contained in a hyperplane; the source defines rigid by requiring its deformations to be induced by automorphisms of $\PP^n$. It does not explicitly require smoothness or irreducibility, fix a field, or distinguish actual deformations from all infinitesimal ones.

The attached aid treats smooth irreducible curves over an algebraically closed field of characteristic zero. For $C\subset\PP^n$, $n\ge1$, write $g=h^1(C,\mathcal O_C)$, $d=\deg C$, and
\[
s=\dim\operatorname{Stab}_{\operatorname{PGL}_{n+1}}(C).
\]
Projective rigidity in the actual Hilbert-deformation sense requires the projective orbit to account for the local reduced deformation locus. The aid states the necessary speciality inequality
\[
(n+1)h^1(C,\mathcal O_C(1))\ge4g-3+s.
\]
Its underlying component-versus-orbit comparison uses the lower bound
\[
\dim_{[C]}\operatorname{Hilb}\ge\chi(N_{C/\PP^n})=(n+1)d+(n-3)(1-g)
\]
and orbit dimension $(n+1)^2-1-s$, together with Riemann--Roch and nondegeneracy. A positive dimension defect supplies actual non-projective Hilbert deformations, not merely surplus potentially obstructed tangent vectors.

The aid gives these consequences: every positive-genus nonspecial embedding, meaning $h^1(C,\mathcal O_C(1))=0$, has a non-projective deformation; no smooth elliptic curve is projectively rigid in any projective space; and among smooth rational nondegenerate curves only rational normal curves are rigid. A rational normal curve is the degree-$n$ image of $\PP^1$ under the complete system $|\mathcal O_{\PP^1}(n)|$. It further reports that these are the only smooth rigid curves for $n\le3$, leaving the special higher-genus case in $n\ge4$ open in its literature assessment. None of these restricted conclusions settles any additional singular or reducible cases allowed by the literal wording.
\subsection{Short English statement}
Ask whether Harris's projective-rigidity phenomenon extends beyond rational normal curves. The supplied aid derives a stabilizer-sensitive speciality obstruction for smooth curves and rules out several classes, while leaving special higher-genus embeddings in higher dimension outside its conclusion.
\subsection{Sources}
\begin{itemize}
\item Local UnsolvedMath registry, permanent ID $20000824$, source code catalogue entry; canonical \texttt{aim-arithmetic-geometry-notes.json}, \texttt{notes[69]}.
\item American Institute of Mathematics, workshop ``Components of Hilbert schemes,'' question attributed to Joe Harris: \url{https://aimath.org/WWN/hilbertschemes/hilbertschemes.pdf}. The metadata gives source item 2, but the extracted number is split as ``2'' and ``3'' and may represent 23. The verification field reports an unambiguous statement on PDF page 2; no independent PDF inspection is claimed here. No original literature or remarks are supplied.
\item Attached research attempt 1: partial result, recorded run status open, rigor proved, novelty confidence low; registry status partially solved, classification PARTIAL-PROGRESS. Its dimension-obstruction and low-dimensional claims are machine generated and not peer reviewed; independently verify before citation.
\end{itemize}
\end{document}
